% ==============================================================================
% PLANTILLA MAESTRA GENERAL - UCNL
% ==============================================================================
% Uso:
% 1. Duplica este archivo y renombra la copia por materia, semana y actividad.
% 2. Actualiza documenttitle, documentsubtitle, documentsubject, coursename,
%    coursecode, universitydepartment y datos de la figura docente.
% 3. Lee la planeacion semanal y NO pegues las instrucciones completas dentro del
%    producto final. Usa la planeacion como lista de cumplimiento.
% 4. Revisa la bibliografia indicada, toma notas, registra citas y redacta con
%    comprension propia.
% 5. Identifica el producto academico solicitado y construyelo dentro del
%    documento siempre que sea posible.
% 6. Entrega en PDF con la nomenclatura indicada en el aula.
%
% Formato institucional recomendado:
% - Fuente tipo Helvetica, equivalente visual a Arial.
% - Interlineado 1.5.
% - Citas autor-anio con natbib: usar \citep{clave} o \citet{clave}.
% - Referencias en bibliografia BibTeX, formato APA 7 en cuerpo y referencias.
% - Caratula, resumen breve, desarrollo, producto solicitado, conclusion y
%   bibliografia.
% - Redaccion academica clara, coherente, cohesionada, formal y sin faltas.
% - La portada puede llevar una marca de agua institucional discreta; ajustar
%   las variables coverwatermark... si se requiere apagarla o cambiar intensidad.
%
% Criterios generales de cumplimiento:
% - Responder exactamente a la actividad solicitada en la planeacion.
% - Identificar conceptos, elementos, criterios o categorias de forma precisa.
% - Analizar, interpretar y tomar postura; no solo copiar fuentes.
% - Organizar la informacion por jerarquia, secuencia y criterios claros.
% - Incluir conclusion personal cuando la actividad lo permita o lo solicite.
% - Usar citas y referencias en APA 7.
% - Evitar copiar las indicaciones de la actividad en el cuerpo final.
%
% Regla para productos visuales:
% - Si la actividad pide tabla, cuadro comparativo, esquema, mapa conceptual,
%   mapa mental, linea de tiempo, diagrama, matriz, lista de verificacion, red
%   conceptual, SQA, organizador grafico o producto similar, realizarlo
%   preferentemente con LaTeX/TikZ para mantener nitidez, evidencia editable y
%   diseno controlado.
% - Si la consigna exige Canva, Miro u otra herramienta externa, insertar una
%   captura nitida y conservar la explicacion dentro del documento.
% - Para tablas extensas, usar pagina limpia, sin encabezados ni pies, landscape,
%   letra pequena y restaurar el estilo al terminar.
%
% Criterios de redaccion personal:
% - No escribir para llenar espacio; escribir para sostener una idea.
% - Convertir informacion en criterio: que significa, como funciona, que problema
%   resuelve, que limite presenta y que consecuencia produce.
% - Estructura mental recomendada: problema -> sistema -> consecuencia -> postura.
% - Usar pensamiento de diagnostico: entrada -> proceso -> salida.
% - Cada parrafo debe tener funcion: presentar, explicar, comparar, valorar o
%   concluir.
% - Usar primera persona con moderacion academica: "Considero que",
%   "Desde esta perspectiva", "La posicion que sostengo es". Evitar "yo creo".
%
% Notas de enfoque personal segun enfoque_personal.txt:
% - La voz debe ser academica, tecnica, directa, estrategica y funcional.
% - El texto no debe sonar como estudiante pasivo que recopila: debe sonar como
%   alguien que interpreta, ordena, cruza datos y busca consecuencia practica.
% - Antes de redactar, identificar la funcion de cada seccion: abrir problema,
%   definir criterio, demostrar con evidencia, comparar, valorar o concluir.
% - Evitar frases genericas como "es importante porque si" o "desde siempre".
%   Sustituirlas por relaciones verificables: causa, efecto, limite, costo,
%   funcion, riesgo, implicacion o mejora.
% - Cada concepto debe pasar de definicion a analisis. Formula util:
%   concepto -> funcion -> riesgo si se aplica mal -> consecuencia -> postura.
% - La postura personal no debe ser opinion suelta. Debe tener tres piezas:
%   posicion, razon y efecto. Ejemplo: "Considero que X, porque Y; por ello Z".
% - Integrar experiencia o criterio propio sin perder formalidad: no narrar la
%   vida personal, sino usar mirada practica para explicar por que el tema
%   importa en educacion, derecho, instituciones, comunicacion o profesion.
%
% Estilo observado en actividades ya realizadas:
% - Abrir con una tesis concreta:
%   "El Derecho no solo ordena conductas: organiza poder, define limites y
%   distribuye consecuencias."
% - Formular el problema en negativo y positivo:
%   "El problema no es solo..., sino..."
% - Convertir el producto visual en argumento:
%   "El cuadro no solo recopila datos; reconstruye la funcion, el problema y el
%   punto que requiere correccion."
% - Cerrar con aprendizaje transferible:
%   "No basta con saber que dice la norma; tambien debe analizarse a quien
%   protege, a quien excluye y que razones permiten defenderla publicamente."
%
% Nivel cognitivo esperado:
% - Recordar: usar solo para definir conceptos basicos. No debe ser el nivel
%   dominante en una actividad ponderable.
% - Comprender: explicar con palabras propias y relacionar con el tema.
% - Aplicar: llevar el concepto a un caso, norma, texto, correo, dilema o hecho.
% - Analizar: separar elementos, detectar causas, efectos, tensiones y limites.
% - Evaluar: sostener una postura con criterios, fuentes y consecuencias.
% - Crear: elaborar producto propio: mapa, cuadro, linea de tiempo, ensayo,
%   propuesta, reformulacion, glosario, portafolio o solucion argumentada.
%
% Regla de nivel minimo:
% - Si la actividad vale puntos y pide producto, no quedarse en recordar o
%   comprender. Debe llegar por lo menos a aplicar y analizar; si pide conclusion
%   o postura, debe llegar a evaluar.
%
% Uso de las 100 tecnicas didacticas:
% - Esta plantilla conserva el manual "Tecnicas didacticas del aprendizaje" que
%   se incorporo en las plantillas de Filosofia del Derecho, Etica y Moral
%   juridica, y Redaccion en contextos virtuales.
% - Cuando la planeacion mencione una tecnica de la coleccion UnADM, identificar:
%   definicion de la tecnica, producto esperado, pasos de construccion, criterios
%   de revision y cierre reflexivo.
% - La tecnica didactica no debe verse como adorno. Debe mostrar seleccion,
%   jerarquia, relacion entre ideas y lectura propia.
% - Siempre que sea viable, construir la tecnica en LaTeX/TikZ: cuadro
%   comparativo, mapa conceptual, esquema, mapa mental, linea de tiempo, matriz,
%   red conceptual, diagrama de flujo, tabla SQA o lista de verificacion.
%
% Checklist antes de entregar:
% [ ] El titulo, materia, semana, docente y fecha estan actualizados.
% [ ] El objetivo coincide con la planeacion.
% [ ] El producto solicitado aparece visible, rotulado y completo.
% [ ] La introduccion plantea un problema o postura, no solo anuncia el tema.
% [ ] El desarrollo explica, conecta y analiza.
% [ ] Hay citas APA autor-anio dentro del texto.
% [ ] La bibliografia final coincide con las fuentes citadas.
% [ ] La conclusion responde que se comprendio y que postura se sostiene.
% [ ] El nivel cognitivo llega al menos a aplicacion y analisis.
% [ ] Si hay conclusion/postura, incluye posicion, razon y consecuencia.
% [ ] La tecnica didactica solicitada esta construida y explicada.
% [ ] No aparecen instrucciones copiadas de la planeacion dentro del producto.
% [ ] Ortografia, acentos, sintaxis y nombres propios revisados.
% [ ] Si se uso IA como apoyo, el texto fue comprendido, revisado y apropiado.
% ==============================================================================

% CREACION DEL DOCUMENTO
\documentclass[
	spanish,
	letterpaper, oneside
]{article}

% ==============================================================================
% INFORMACION DEL DOCUMENTO
% Actualizar en cada actividad.
% Ejemplos:
% \def\documenttitle {Cuadro comparativo de conceptos fundamentales}
% \def\documentsubtitle {Actividad 1 - Matematicas I}
% \def\documentsubject {Actividad 1}
% ==============================================================================
\def\documenttitle {Actividad 1 - Matematicas I}
\def\documentsubtitle {Actividad 1 - Matematicas I}
\def\documentsubject {Actividad 1}

\def\documentauthor {Martin Jonathan de la Cruz}
\def\coursename {Matematicas I}
\def\coursecode {MAT-I}

\def\universityname {Universidad Ciudadana de Nuevo Leon}
\def\universityfaculty {Programa academico UCNL}
\def\universitydepartment {Matematicas I}
\def\universitydepartmentimage {departamentos/logo-ucnl}
\def\universitydepartmentimagecfg {height=1.57cm}
\def\universitylocation {Monterrey, Nuevo Leon}

% INTEGRANTES, PROFESORES Y FECHAS
\def\authortable {
	\begin{tabular}{ll}
		\textbf{Alumno:}
		& \begin{tabular}[t]{l}
			 Martin Jonathan de la Cruz \\
		\end{tabular} \\
		Matricula: & UCNL \\

		\textit{Figura docente:}
		& \begin{tabular}[t]{l}
			Nombre \\ Apellido
		\end{tabular} \\
		& \\
		\multicolumn{2}{l}{Fecha de realizacion: \today} \\
		\multicolumn{2}{l}{\universitylocation}
	\end{tabular}
}

% IMPORTACION DEL TEMPLATE BASE
\input{template}

% FORMATO GENERAL
\usepackage{helvet}
\renewcommand{\familydefault}{\sfdefault}

% TABLAS Y PRODUCTOS VISUALES
\usepackage{array}
\usepackage{booktabs}
\usepackage{longtable}
\usepackage{pdflscape}
\usepackage{tikz}
\usetikzlibrary{
	arrows.meta,
	positioning,
	calc,
	matrix,
	fit,
	shapes.geometric,
	shadows.blur
}

% CITAS EN FORMATO AUTOR-ANIO, COMPATIBLE CON APA 7 EN EL CUERPO DEL TEXTO
% Usar:
% - Cita parentetica: \citep{clave}
% - Cita narrativa: \citet{clave}
% - Varias fuentes: \citep{clave1,clave2}
\setcitestyle{authoryear,open={(},close={)}}

% ==============================================================================
% AYUDAS DE REDACCION
% Quitar o comentar \pendiente cuando el documento final este terminado.
% ==============================================================================
\newcommand{\pendiente}[1]{}

% ==============================================================================
% MARCA DE AGUA EN PORTADA
% - Se aplica solo a la portada porque usa \AddToShipoutPictureBG*.
% - Para apagarla en alguna actividad: \def\coverwatermarkenabled {false}
% - Usar opacidad baja para que no compita con titulo, datos ni logotipo.
% ==============================================================================
\def\coverwatermarkenabled {true}
\def\coverwatermarkimage {img/departamentos/logo-nuevo-leon.png}
\def\coverwatermarkopacity {0.16}
\def\coverwatermarkwidth {0.72\paperwidth}
\def\coverwatermarkxshift {0cm}
\def\coverwatermarkyshift {-0.35cm}

\newcommand{\insertcoverwatermark}{%
	\ifthenelse{\equal{\coverwatermarkenabled}{true}}{%
		\AddToShipoutPictureBG*{%
			\begin{tikzpicture}[remember picture,overlay]
				\node[
					opacity=\coverwatermarkopacity,
					inner sep=0pt,
					xshift=\coverwatermarkxshift,
					yshift=\coverwatermarkyshift
				] at (current page.center) {%
					\includegraphics[width=\coverwatermarkwidth]{\coverwatermarkimage}%
				};
			\end{tikzpicture}%
		}%
	}{}%
}

\begin{document}

% PORTADA
\insertcoverwatermark
\templatePortrait

% CONFIGURACION DE PAGINA Y ENCABEZADOS
\templatePagecfg
\onehalfspacing

% ==============================================================================
% RESUMEN / ABSTRACT
% Debe ser breve. Formula recomendada:
% Tema + objetivo + tecnica/producto + enfoque personal.
% ==============================================================================
\begin{abstractd}
	\textbf{Objetivo:} Resolver y clasificar expresiones algebraicas basicas mediante ecuaciones lineales, ecuaciones cuadraticas, desigualdades, valor absoluto, dominio de funciones y composicion de funciones.

	\textbf{Producto elaborado:} Cuestionario resuelto de opcion multiple sobre contenidos introductorios de Matematicas I.

	\textbf{Enfoque de analisis:} Aplicar procedimientos directos de despeje, sustitucion y verificacion para identificar la respuesta correcta en cada reactivo y distinguir entre operaciones validas, restricciones de dominio y tipos de expresion algebraica.

	\textbf{Nivel cognitivo:} Aplicar y analizar, porque cada reactivo exige interpretar la expresion dada, ejecutar el procedimiento correspondiente y validar la solucion obtenida.
\end{abstractd}

% TABLA DE CONTENIDOS - INDICE
\templateIndex

% CONFIGURACIONES FINALES
\templateFinalcfg

% ======================= INICIO DEL DOCUMENTO =======================

% ==============================================================================
% SECCION 1: INTRODUCCION
% Apertura recomendada:
% - Iniciar con una afirmacion clara, no con una frase generica.
% - Traducir la consigna a un problema academico, juridico, etico, comunicativo
%   o profesional, segun la materia.
% - Explicar por que importa.
% - Cerrar anticipando el enfoque personal.
% ==============================================================================
\section{Introduccion}

El estudio inicial del algebra permite convertir situaciones simbolicas en procedimientos ordenados de resolucion. En Matematicas I, esta base resulta indispensable porque conecta el significado de variable, constante, igualdad, desigualdad, dominio y funcion con operaciones concretas de despeje, evaluacion y verificacion.

La actividad atiende el problema de reconocer que no basta con memorizar definiciones: es necesario identificar que tipo de expresion se tiene enfrente y aplicar la regla adecuada sin perder de vista sus restricciones. Un error pequeno en signos, sustitucion o interpretacion del dominio modifica por completo el resultado.

Desde esta perspectiva, el cuestionario se resuelve con un criterio operativo: en cada reactivo se determina primero la estructura matematica, despues se ejecuta el procedimiento minimo necesario y, por ultimo, se contrasta la respuesta con las opciones disponibles.

% ==============================================================================
% SECCION 2: DESARROLLO / ANALISIS
% Adaptar los subtitulos segun la materia y la planeacion.
%
% Productos posibles:
% - Cuadro comparativo: criterios, evidencia, diferencia, valoracion.
% - Mapa conceptual: concepto central, conceptos secundarios, conectores claros.
% - Linea de tiempo: fecha/etapa, hecho, aporte, consecuencia.
% - Estudio de caso: situacion, elementos, problema, interpretacion, solucion.
% - Lista de verificacion: criterio observable, indicador, cumplimiento, mejora.
% - Glosario: concepto, definicion, funcion, ejemplo, fuente.
% - Reporte o ensayo: tesis, argumentos, evidencia, postura y cierre.
%
% Regla de parrafo:
% Entrada: concepto o problema.
% Proceso: explicacion, fuente y relacion.
% Salida: consecuencia, criterio o postura.
% ==============================================================================
\section{Desarrollo}

\subsection{Contexto y problema}

El cuestionario se ubica en la etapa introductoria del curso, donde el estudiante debe demostrar dominio sobre conceptos elementales del lenguaje algebraico. El problema de analisis consiste en distinguir entre ecuaciones, desigualdades y funciones para resolver cada caso con el metodo correcto y no por simple intuicion.

\subsection{Marco conceptual}

Los conceptos centrales de la actividad son variable, constante, coeficiente, ecuacion lineal, ecuacion cuadratica, valor absoluto, dominio y composicion de funciones. Cada uno cumple una funcion precisa: la variable representa un valor que puede cambiar, la constante fija un dato numerico, la ecuacion expresa una igualdad que debe satisfacerse y el dominio delimita los valores para los que una funcion tiene sentido. Comprender estas nociones evita errores frecuentes, como dividir entre cero, confundir una constante con una incognita o ignorar que el valor absoluto genera soluciones simetricas.

\subsection{Nivel cognitivo aplicado}

La actividad exige principalmente aplicar procedimientos algebraicos y analizar la forma de cada expresion. Resolver no implica solo obtener un numero: tambien requiere clasificar el tipo de ecuacion, interpretar el signo de desigualdad, evaluar funciones en un punto y revisar si una opcion de respuesta realmente coincide con el resultado calculado.

\subsection{Producto solicitado}

Se presenta el cuestionario resuelto de la Actividad 1.

\clearpage
\pagestyle{empty}
\begin{landscape}
\begingroup
\scriptsize
\setlength{\tabcolsep}{4pt}\renewcommand{\arraystretch}{1.25}
\begin{longtable}{@{}>{\centering\arraybackslash}p{0.035\linewidth}>{\raggedright\arraybackslash}p{0.42\linewidth}>{\raggedright\arraybackslash}p{0.20\linewidth}>{\raggedright\arraybackslash}p{0.30\linewidth}@{}}
\caption{Cuestionario resuelto de conceptos introductorios de Matematicas I.}\label{tab:reporte-matematicas-i-actividad-1}\\
\toprule
\textbf{No.} & \textbf{Pregunta y opciones} & \textbf{Respuesta correcta} & \textbf{Justificación} \\
\midrule
\endfirsthead
\toprule
\textbf{No.} & \textbf{Pregunta y opciones} & \textbf{Respuesta correcta} & \textbf{Justificación} \\
\midrule
\endhead
\bottomrule
\endfoot
\bottomrule
\endlastfoot
1 & Resolver $4x=-20$. \newline a) $x=-4$   b) $x=-5$   c) $x=5$   d) $x=20$. & \textbf{b) $x=-5$.} & Es correcta porque $x=-5$ corresponde con precisión al concepto evaluado en el reactivo; las demás opciones no describen adecuadamente esa relación. \\
\midrule
2 & En la ecuacion $x+2=3$, el numero 2 es... \newline a) una raiz   b) una constante   c) un coeficiente cuadratico   d) una variable. & \textbf{b) una constante.} & Es correcta porque una constante corresponde con precisión al concepto evaluado en el reactivo; las demás opciones no describen adecuadamente esa relación. \\
\midrule
3 & ?Que tipo de ecuacion es $ax+b=0$? \newline a) trigonometrica   b) cuadratica   c) cubica   d) lineal. & \textbf{d) lineal.} & Es correcta porque lineal corresponde con precisión al concepto evaluado en el reactivo; las demás opciones no describen adecuadamente esa relación. \\
\midrule
4 & Resolver $5-x=-10$. \newline a) $x=-15$   b) $x=15$   c) $x=-10$   d) $x=10$. & \textbf{b) $x=15$.} & Es correcta porque $x=15$ corresponde con precisión al concepto evaluado en el reactivo; las demás opciones no describen adecuadamente esa relación. \\
\midrule
5 & Evaluar $F(x)=5x-2$ en $x=-1$. \newline a) $-7$   b) $5$   c) $7$   d) $-5$. & \textbf{a) $-7$.} & Es correcta porque $-7$ corresponde con precisión al concepto evaluado en el reactivo; las demás opciones no describen adecuadamente esa relación. \\
\midrule
6 & ?Que garantiza que dos ecuaciones sean equivalentes? \newline a) que tengan el mismo numero de terminos   b) que ambas sean lineales   c) que se parezcan   d) que tengan exactamente las mismas soluciones. & \textbf{d) que tengan exactamente las mismas soluciones.} & Es correcta porque que tengan exactamente las mismas soluciones corresponde con precisión al concepto evaluado en el reactivo; las demás opciones no describen adecuadamente esa relación. \\
\midrule
7 & ?Que representa una desigualdad? \newline a) una igualdad perfecta   b) una comparacion entre cantidades donde no necesariamente son iguales   c) una ecuacion cuadratica   d) una funcion racional. & \textbf{b) una comparacion entre cantidades donde no necesariamente son iguales.} & Es correcta porque una comparacion entre cantidades donde no necesariamente son iguales corresponde con precisión al concepto evaluado en el reactivo; las demás opciones no describen adecuadamente esa relación. \\
\midrule
8 & El dominio de $F(x)=\frac{1}{x}$ excluye... \newline a) ningun numero   b) $x=0$   c) $x=-1$   d) $x=1$. & \textbf{b) $x=0$.} & Es correcta porque $x=0$ corresponde con precisión al concepto evaluado en el reactivo; las demás opciones no describen adecuadamente esa relación. \\
\midrule
9 & Resolver $|2x|=8$. \newline a) $x=2$   b) $x=-8$   c) $x=4$ o $x=-4$   d) $x=8$. & \textbf{c) $x=4$ o $x=-4$.} & Es correcta porque $x=4$ o $x=-4$ corresponde con precisión al concepto evaluado en el reactivo; las demás opciones no describen adecuadamente esa relación. \\
\midrule
10 & ?Cual no es una caracteristica de una variable? \newline a) representa valores desconocidos   b) puede cambiar de valor   c) es siempre un numero fijo   d) puede ser $x$, $y$ o $z$. & \textbf{c) es siempre un numero fijo.} & Es correcta porque es siempre un numero fijo corresponde con precisión al concepto evaluado en el reactivo; las demás opciones no describen adecuadamente esa relación. \\
\midrule
11 & Si $F(x)=x+1$ y $G(x)=x^2$, entonces $F(G(x))=$ \newline a) $x^2-1$   b) $x+x^2$   c) $x^2+1$   d) $x^3$. & \textbf{c) $x^2+1$.} & Es correcta porque $x^2+1$ corresponde con precisión al concepto evaluado en el reactivo; las demás opciones no describen adecuadamente esa relación. \\
\midrule
12 & Resolver $x^2-1=0$. \newline a) $x=0$   b) $x=2$   c) $x=1$   d) $x=\pm 1$. & \textbf{d) $x=\pm 1$.} & Es correcta porque $x=\pm 1$ corresponde con precisión al concepto evaluado en el reactivo; las demás opciones no describen adecuadamente esa relación. \\
\midrule
13 & Si $F(x)=\frac{3x+1}{5x-8}$, entonces $F(2)$... \newline Opciones dadas: a) es 0   b) es 2   c) es 1   d) no se puede dividir porque no es entero para ese valor.\\
	Observacion: al sustituir $x=2$, se obtiene $F(2)=\frac{7}{2}$, por lo que ninguna opcion coincide exactamente con el resultado correcto. & \textbf{Observacion:.} & Es correcta porque observacion: corresponde con precisión al concepto evaluado en el reactivo; las demás opciones no describen adecuadamente esa relación. \\
\midrule
14 & El valor absoluto de un numero representa... \newline a) la mitad del numero   b) el signo contrario del numero   c) la distancia del numero al cero en la recta numerica   d) el doble del numero. & \textbf{c) la distancia del numero al cero en la recta numerica.} & Es correcta porque la distancia del numero al cero en la recta numerica corresponde con precisión al concepto evaluado en el reactivo; las demás opciones no describen adecuadamente esa relación. \\
\midrule
15 & Determinar $F(0)$ en $F(x)=4x+1$. \newline a) $1$   b) $4$   c) $-1$   d) $0$. & \textbf{a) $1$.} & Es correcta porque $1$ corresponde con precisión al concepto evaluado en el reactivo; las demás opciones no describen adecuadamente esa relación. \\
\midrule
16 & ?Como se llaman los numeros fijos en una ecuacion? \newline a) coeficientes   b) constantes   c) variables   d) incognitas. & \textbf{b) constantes.} & Es correcta porque constantes corresponde con precisión al concepto evaluado en el reactivo; las demás opciones no describen adecuadamente esa relación. \\
\midrule
17 & Resolver $x^2-4=0$. \newline a) $x=0$   b) $x=\pm 2$   c) $x=4$   d) $x=-4$. & \textbf{b) $x=\pm 2$.} & Es correcta porque $x=\pm 2$ corresponde con precisión al concepto evaluado en el reactivo; las demás opciones no describen adecuadamente esa relación. \\
\midrule
18 & Resolver $x+12=20$. \newline a) $x=20$   b) $x=8$   c) $x=-20$   d) $x=-8$. & \textbf{b) $x=8$.} & Es correcta porque $x=8$ corresponde con precisión al concepto evaluado en el reactivo; las demás opciones no describen adecuadamente esa relación. \\
\midrule
19 & El dominio natural de una funcion es... \newline a) los numeros enteros unicamente   b) los valores positivos   c) el mayor conjunto de numeros reales donde la regla tiene sentido   d) el conjunto solucion. & \textbf{c) el mayor conjunto de numeros reales donde la regla tiene sentido.} & Es correcta porque el mayor conjunto de numeros reales donde la regla tiene sentido corresponde con precisión al concepto evaluado en el reactivo; las demás opciones no describen adecuadamente esa relación. \\
\midrule
20 & $|x|=5$ tiene como solucion... \newline a) $x=-5$ unicamente   b) $x=0$   c) $x=5$ y $x=-5$   d) $x=5$ unicamente. & \textbf{c) $x=5$ y $x=-5$.} & Es correcta porque $x=5$ y $x=-5$ corresponde con precisión al concepto evaluado en el reactivo; las demás opciones no describen adecuadamente esa relación. \\
\midrule
21 & Resolver $x^2-121=0$. \newline a) $x=\pm 11$   b) $x=-11$   c) $x=0$   d) $x=11$. & \textbf{a) $x=\pm 11$.} & Es correcta porque $x=\pm 11$ corresponde con precisión al concepto evaluado en el reactivo; las demás opciones no describen adecuadamente esa relación. \\
\midrule
22 & La composicion $F(G(x))$ significa... \newline a) multiplicar funciones   b) dividir funciones   c) sumar funciones   d) sustituir $G(x)$ dentro de $F(x)$. & \textbf{d) sustituir $G(x)$ dentro de $F(x)$.} & Es correcta porque sustituir $G(x)$ dentro de $F(x)$ corresponde con precisión al concepto evaluado en el reactivo; las demás opciones no describen adecuadamente esa relación. \\
\midrule
23 & Una ecuacion lineal puede representarse como... \newline a) $ax^3+bx+c=0$   b) $mx+b=0$   c) $\frac{a}{x}=b$   d) $ax^2+bx+c=0$. & \textbf{b) $mx+b=0$.} & Es correcta porque $mx+b=0$ corresponde con precisión al concepto evaluado en el reactivo; las demás opciones no describen adecuadamente esa relación. \\
\midrule
24 & ?Que indica el signo $=$ en una ecuacion? \newline a) que ambas expresiones tienen el mismo valor   b) que la ecuacion no tiene solucion   c) que se multiplica   d) que se divide. & \textbf{a) que ambas expresiones tienen el mismo valor.} & Es correcta porque que ambas expresiones tienen el mismo valor corresponde con precisión al concepto evaluado en el reactivo; las demás opciones no describen adecuadamente esa relación. \\
\midrule
25 & Resolver $|x+2|>3$. \newline a) $-3<x<3$   b) $x<-5$ o $x>1$   c) $x=3$   d) $x=-3$. & \textbf{b) $x<-5$ o $x>1$.} & Es correcta porque $x<-5$ o $x>1$ corresponde con precisión al concepto evaluado en el reactivo; las demás opciones no describen adecuadamente esa relación. \\
\midrule
26 & El dominio de $F(x)=\sqrt{x}$ es... \newline a) $x\leq 0$   b) $x\geq 0$   c) todos los reales   d) solo enteros. & \textbf{b) $x\geq 0$.} & Es correcta porque $x\geq 0$ corresponde con precisión al concepto evaluado en el reactivo; las demás opciones no describen adecuadamente esa relación. \\
\midrule
27 & La ecuacion $x^2-9=0$ es de tipo... \newline a) exponencial   b) lineal   c) fraccionaria   d) cuadratica. & \textbf{d) cuadratica.} & Es correcta porque cuadratica corresponde con precisión al concepto evaluado en el reactivo; las demás opciones no describen adecuadamente esa relación. \\
\midrule
28 & Resolver $|x|\geq 6$. \newline a) $-6<x<6$   b) $x=6$   c) $x\leq -6$ o $x\geq 6$   d) $-6\leq x\leq 6$. & \textbf{c) $x\leq -6$ o $x\geq 6$.} & Es correcta porque $x\leq -6$ o $x\geq 6$ corresponde con precisión al concepto evaluado en el reactivo; las demás opciones no describen adecuadamente esa relación. \\
\midrule
29 & Si $F(x)=\frac{3x}{x-2}$, entonces $x\neq$ \newline a) $2$   b) $1$   c) $0$   d) $-2$. & \textbf{a) $2$.} & Es correcta porque $2$ corresponde con precisión al concepto evaluado en el reactivo; las demás opciones no describen adecuadamente esa relación. \\
\midrule
30 & Resolver $7x=49$. \newline a) $x=14$   b) $x=7$   c) $x=-7$   d) $x=1$. & \textbf{b) $x=7$.} & Es correcta porque $x=7$ corresponde con precisión al concepto evaluado en el reactivo; las demás opciones no describen adecuadamente esa relación. \\
\end{longtable}
\endgroup
\end{landscape}
\pagestyle{fancy}
\clearpage

\subsection{Tecnica didactica aplicada}

En este desarrollo se utilizo la tecnica de cuestionario estructurado, porque permite verificar de manera puntual la comprension de conceptos y procedimientos basicos. La secuencia de reactivos favorece la identificacion rapida del tipo de problema, la aplicacion del algoritmo correspondiente y la comprobacion inmediata del resultado.

% ==============================================================================
% BLOQUE OPCIONAL: TABLA O CUADRO COMPARATIVO AMPLIO
% Usar cuando el producto necesite mucho ancho:
% - \clearpage para iniciar en pagina limpia.
% - \pagestyle{empty} para quitar encabezado y pie en el producto amplio.
% - landscape para aprovechar el ancho horizontal.
% - scriptsize, tabcolsep y arraystretch para controlar legibilidad.
% - Restaurar \pagestyle{fancy} al terminar.
%
% Si la tabla es corta, puede omitirse landscape y usarse longtable normal.
% ==============================================================================
% \clearpage
% \pagestyle{empty}
% \begin{landscape}
% \begingroup
% \scriptsize
% \setlength{\tabcolsep}{4pt}
% \renewcommand{\arraystretch}{1.35}
% \begin{longtable}{@{}p{0.18\linewidth}p{0.39\linewidth}p{0.39\linewidth}@{}}
% \caption{Titulo del cuadro comparativo}\label{tab:cuadro-comparativo}\\
% \toprule
% \textbf{Criterio} & \textbf{Elemento 1} & \textbf{Elemento 2} \\
% \midrule
% \endfirsthead
% \toprule
% \textbf{Criterio} & \textbf{Elemento 1} & \textbf{Elemento 2} \\
% \midrule
% \endhead
% Concepto & \textit{Refuerzo Ciclo A: Contenido sintetico con cita si aplica.} & \textit{Refuerzo Ciclo A: Contenido sintetico con cita si aplica.} \\
% Caracteristicas & \textit{Refuerzo Ciclo A: Contenido.} & \textit{Refuerzo Ciclo A: Contenido.} \\
% Funcion & \textit{Refuerzo Ciclo A: Contenido.} & \textit{Refuerzo Ciclo A: Contenido.} \\
% Valoracion critica & \textit{Refuerzo Ciclo A: Contenido.} & \textit{Refuerzo Ciclo A: Contenido.} \\
% \bottomrule
% \end{longtable}
% \endgroup
% \end{landscape}
% \pagestyle{fancy}

% ==============================================================================
% BLOQUE OPCIONAL: PRODUCTO VISUAL EN TIKZ
% Usar para mapa conceptual, mapa mental, esquema, flujo, proceso o diagrama.
% Cada flecha debe expresar relacion; no usar nodos sueltos sin conectores.
% ==============================================================================
% \begin{figure}[H]
% 	\centering
% 	\begin{tikzpicture}[
% 		node distance=1.3cm,
% 		box/.style={rectangle, rounded corners=2pt, draw=black!65, fill=black!4, align=center, minimum width=3.2cm, minimum height=0.9cm},
% 		arrow/.style={-{Latex[length=2.5mm]}, thick, draw=black!65}
% 	]
% 		\node[box] (centro) {Concepto central};
% 		\node[box, right=of centro] (relacion) {Concepto relacionado};
% 		\node[box, below=of relacion] (efecto) {Consecuencia};
%
% 		\draw[arrow] (centro) -- node[above]{implica} (relacion);
% 		\draw[arrow] (relacion) -- node[right]{produce} (efecto);
% 	\end{tikzpicture}
% 	\caption{Titulo del producto visual.}
% \end{figure}

\subsection{Analisis propio}

El cuestionario muestra que los contenidos iniciales de algebra no deben estudiarse como temas aislados. Ecuaciones, desigualdades, valor absoluto y dominio de funciones comparten una misma exigencia: interpretar correctamente la expresion antes de operar. El principal hallazgo es que muchos reactivos se resuelven con procedimientos breves, pero solo si se reconoce con precision la estructura matematica involucrada. Por ello, el aprendizaje relevante no es solo obtener la opcion correcta, sino justificar por que las demas opciones no satisfacen la condicion planteada.

% ==============================================================================
% SECCION 3: POSTURA PERSONAL
% Debe mostrar comprension real. No repetir fuentes; interpretar.
% ==============================================================================
\section{Postura personal}

Considero que este tipo de cuestionario cumple una funcion formativa importante en Matematicas I, porque obliga a consolidar operaciones basicas que despues se vuelven necesarias en problemas mas complejos. La razon es que el algebra inicial entrena orden, precision y verificacion. La consecuencia de dominar estos fundamentos es que el estudiante puede avanzar hacia contenidos posteriores con menor riesgo de arrastrar errores conceptuales.

% ==============================================================================
% SECCION 4: CONCLUSION
% Cierre recomendado:
% - Responder directamente a la actividad.
% - Indicar que se comprendio, no solo que se elaboro.
% - Mencionar criterio personal final.
% - No introducir informacion nueva.
% ==============================================================================

\section{Conclusion}

El desarrollo de la Actividad 1 permitio resolver reactivos relacionados con ecuaciones lineales y cuadraticas, desigualdades, valor absoluto, dominio y composicion de funciones. La revision muestra que cada resultado correcto depende de interpretar bien la expresion y aplicar un procedimiento coherente con su estructura. En conclusion, el cuestionario no solo evalua memoria operativa, sino comprension matematica elemental, y por ello constituye una base util para el resto del curso.

% ==============================================================================
% MANUAL DE TECNICAS DIDACTICAS
% Dejar incluido en la plantilla maestra. En una entrega final, puede comentarse
% esta linea si la docente solo solicita el producto de la semana.
% ==============================================================================
% Referencia opcional desactivada: archivo externo no disponible en este repositorio.

% BIBLIOGRAFIA
% Cambiar el nombre del archivo .bib segun la materia/actividad.
% Ejemplo: \bibliography{filosofia-del-derecho}
\bibliography{matematicas-I}

% FIN DEL DOCUMENTO





































































































\end{document}


